\documentclass{article}
\usepackage{amsmath,amssymb}
\title{Why higher feedback evaluation needs higher uniformity}
\date{15 November 2026}
\begin{document}\maketitle
\paragraph{Status.} A counterboundary for the accompanying conditional evaluation construction. It tests a proposed extension of the earlier two-dimensional evaluation doctrine to bicategorical states.

\paragraph{Claim.} An evaluation specified on state objects, state maps and their ordinary directed squares need not extend to the full native state tricategory while preserving that lower-dimensional assignment. A target must also accommodate the admitted state two-cells and their native three-cell equations.

\paragraph{Test case.} Fix universes for small categories. Take the symmetric monoidal state bicategory $S=(\mathbf{Cat},\times)$ and the terminal process bicategory $K_0$, with its constant pure state action. There is one external interface and every process body and process square is unique. The resulting normalized native state tricategory is the one-object tricategory $BS$: its one-cells are state categories, its two-cells are functors, and its three-cells are natural transformations. Composition is the chosen Cartesian product with its actual coherence.

Let $S_1$ be the ordinary monoidal category of the same small categories and functors, obtained by discarding natural transformations. This describes an underlying ordinary category, not an asserted higher forgetful functor. Its one-object delooping $B S_1$ is a bicategory, regarded as a tricategory with only identity three-cells. The ordinary category-state construction over the same terminal process doctrine gives precisely this target. Thus the identity assignment on state categories and functors is a valid comparison for the lower-dimensional state presentation.

Take $X$ to be the two-object poset category $0\to1$. The endofunctors $h$ and $k$ constantly equal to $0$ and $1$ are distinct, but there is a natural transformation $\sigma:h\Rightarrow k$. In the source native tricategory it is an admitted three-cell: the process admission equation is automatic because the process doctrine is terminal. The target has no three-cell from $h$ to $k$, since these are distinct target two-cells. Consequently there is no trihomomorphism extending the specified identity lower-dimensional assignment. This is a cell-existence obstruction, independent of choices of associators or unitors.

\paragraph{Interpretation.} The example does not rule out feedback evaluations that collapse more data, nor evaluations into a target with appropriate nontrivial three-cells. It rules out inferring the higher evaluator from an evaluator of the ordinary state-map presentation alone. A general native evaluation theorem must explicitly specify the image of every admitted three-cell, its functoriality, and its interaction with sequential/parallel assembly and joining. The generated tower construction supplies source coherence, not those additional semantic target data. No instantaneous feedback obstruction or new state semantics is claimed.
\end{document}
